\documentclass[11pt,a4paper]{article}
\usepackage{amsmath,amsthm,amssymb,amsfonts}
\usepackage{mathtools}
\usepackage{hyperref}
\usepackage{cleveref}
\usepackage{enumitem}
\usepackage{tikz-cd}
\usepackage{booktabs}
\usepackage{xcolor}
\usepackage{tcolorbox}
\usepackage{geometry}
\geometry{margin=1in}

% Theorem environments
\newtheorem{theorem}{Theorem}[section]
\newtheorem{lemma}[theorem]{Lemma}
\newtheorem{proposition}[theorem]{Proposition}
\newtheorem{corollary}[theorem]{Corollary}
\newtheorem{conjecture}[theorem]{Conjecture}
\theoremstyle{definition}
\newtheorem{definition}[theorem]{Definition}
\newtheorem{example}[theorem]{Example}
\newtheorem{remark}[theorem]{Remark}
\newtheorem{assumption}{Assumption}

% Custom commands
\newcommand{\C}{\mathsf{C}}
\newcommand{\Cstar}{\mathsf{C}^*}
\newcommand{\tmin}{\otimes_{\min}}
\newcommand{\tmax}{\otimes_{\max}}
\newcommand{\R}{\mathbb{R}}
\newcommand{\id}{\mathrm{id}}
\newcommand{\interior}{\mathrm{int}}

% Status boxes
\newtcolorbox{statusbox}[2][]{colback=#2!5!white,colframe=#2!75!black,title=#1}
\newtcolorbox{challengebox}{colback=red!5!white,colframe=red!75!black,title=Challenge}
\newtcolorbox{rejectedbox}{colback=red!10!white,colframe=red!90!black,title=REJECTED}

\title{Proof Attempt: GJWH Generalization for Proper Cones\\[0.5em]
\large An Alethfeld Orchestrator Documentation}
\author{Generated by Alethfeld Proof Orchestrator v5.1}
\date{January 11, 2026\\Graph ID: \texttt{graph-aa3f7c-d06137}, Version 15}

\begin{document}
\maketitle

\begin{abstract}
This document provides a comprehensive record of an attempted proof of the GJWH (Gisin--Hughston--Jozsa--Wootters) theorem generalized from positive semidefinite cones to arbitrary proper convex cones. The proof attempt followed an explicit construction strategy recommended by the Adviser subagent. Six parallel Verifier subagents identified critical gaps, with three proof steps ultimately rejected. The fundamental obstruction discovered is that the explicit construction can only produce scalar rescalings of fixed directions, whereas genuine steering assemblages require ``rotations'' within the cone that the maximal tensor product cannot support for general proper cones.
\end{abstract}

\tableofcontents
\newpage

%==============================================================================
\section{Introduction}
%==============================================================================

\subsection{The Classical GJWH Theorem}

The Gisin--Hughston--Jozsa--Wootters (GJWH) theorem is a fundamental result in quantum information theory. In its original formulation for quantum states (positive semidefinite matrices), it states:

\begin{theorem}[GJWH for Quantum States]
Let $\{\rho_{a|x}\}_{a,x}$ be a steering assemblage on a finite-dimensional Hilbert space $\mathcal{H}_B$, meaning:
\begin{enumerate}[label=(\roman*)]
    \item Each $\rho_{a|x}$ is a positive semidefinite operator on $\mathcal{H}_B$
    \item $\sum_a \rho_{a|x} = \rho_B$ for all measurement settings $x$ (no-signaling)
    \item $\mathrm{Tr}(\rho_B) = 1$ (normalization)
\end{enumerate}
Then there exists a bipartite state $\rho_{AB}$ on $\mathcal{H}_A \otimes \mathcal{H}_B$ and POVMs $\{M_{a|x}\}$ on $\mathcal{H}_A$ such that:
\[
\rho_{a|x} = \mathrm{Tr}_A\left[(M_{a|x} \otimes I_B)\rho_{AB}\right]
\]
\end{theorem}

This theorem establishes that any steering assemblage can be ``explained'' by measurements on one half of a shared bipartite quantum state.

\subsection{The Proposed Generalization}

The theorem we attempted to prove generalizes this result from positive semidefinite cones to arbitrary proper convex cones:

\begin{theorem}[Claimed GJWH Generalization]\label{thm:main}
Let $V$ be a finite-dimensional real vector space and $\C \subset V$ a proper cone (closed, pointed, generating). Let $\phi \in \interior(\Cstar)$ be a linear functional in the interior of the dual cone. If $\{v_{a|x}\}_{a,x}$ is a steering assemblage in $\C$, meaning:
\begin{enumerate}[label=(\roman*)]
    \item $v_{a|x} \in \C$ for all outcomes $a$ and settings $x$
    \item $\sum_a v_{a|x} = v_*$ is independent of $x$ (no-signaling)
    \item $\phi(v_*) = 1$ (normalization)
\end{enumerate}
Then there exists a set of measurements $f_{a|x} \in \Cstar$ with $\sum_a f_{a|x} = \phi$, and a vector $w \in \C \tmax \C$, such that:
\[
v_{a|x} = (f_{a|x} \otimes \id)(w)
\]
\end{theorem}

\subsection{Proof Mode and Methodology}

This proof attempt was conducted in \textbf{strict-mathematics} mode using the Alethfeld Proof Orchestrator v5.1. The methodology involved:

\begin{enumerate}
    \item \textbf{Theorem Audit}: Preliminary plausibility assessment
    \item \textbf{Strategy Evaluation}: Selection of proof approach
    \item \textbf{Skeleton Generation}: Top-level proof structure
    \item \textbf{Parallel Verification}: Six independent Verifier subagents
    \item \textbf{Gap Diagnosis}: Adviser analysis of discovered obstructions
\end{enumerate}

%==============================================================================
\section{Preliminary Definitions}
%==============================================================================

\begin{definition}[Proper Cone]
A subset $\C \subset V$ of a finite-dimensional real vector space is a \emph{proper cone} if it satisfies:
\begin{enumerate}[label=(\roman*)]
    \item \textbf{Cone property}: $\lambda x \in \C$ for all $x \in \C$ and $\lambda \geq 0$
    \item \textbf{Convexity}: $x + y \in \C$ for all $x, y \in \C$
    \item \textbf{Closed}: $\C$ is topologically closed
    \item \textbf{Pointed}: $\C \cap (-\C) = \{0\}$
    \item \textbf{Generating}: $\C - \C = V$ (equivalently, $\C$ has nonempty interior)
\end{enumerate}
\end{definition}

\begin{definition}[Dual Cone]
The \emph{dual cone} of $\C \subset V$ is:
\[
\Cstar = \{f \in V^* : f(x) \geq 0 \text{ for all } x \in \C\}
\]
For a proper cone, $\Cstar$ is also proper, and $(\Cstar)^* = \C$ by the bipolar theorem.
\end{definition}

\begin{definition}[Tensor Products of Cones]
For proper cones $\C_A \subset V_A$ and $\C_B \subset V_B$:
\begin{align}
\C_A \tmin \C_B &= \overline{\mathrm{conv}}\{x \otimes y : x \in \C_A, y \in \C_B\} \\
\C_A \tmax \C_B &= \{z \in V_A \otimes V_B : (f \otimes g)(z) \geq 0 \text{ for all } f \in \C_A^*, g \in \C_B^*\}
\end{align}
The minimal tensor product consists of ``separable'' elements; the maximal contains additional ``entangled'' elements satisfying all bilinear positivity constraints.
\end{definition}

%==============================================================================
\section{Assumptions (Depth 0 Nodes)}
%==============================================================================

The proof begins with four explicit assumptions:

\begin{assumption}[A1: \texttt{:0-a00001}]
$V$ is a finite-dimensional real vector space.
\end{assumption}

\begin{assumption}[A2: \texttt{:0-a00002}]
$\C \subset V$ is a proper cone (closed, pointed, generating).
\end{assumption}

\begin{assumption}[A3: \texttt{:0-a00003}]
$\phi \in \interior(\Cstar)$.
\end{assumption}

\begin{assumption}[A4: \texttt{:0-a00004}]
$\{v_{a|x}\}_{a,x}$ is a steering assemblage in $\C$: each $v_{a|x} \in \C$, the sum $\sum_a v_{a|x} = v_*$ is independent of $x$, and $\phi(v_*) = 1$.
\end{assumption}

\begin{remark}
These assumptions are standard for the generalization. Note that A3 ($\phi \in \interior(\Cstar)$) ensures that $\phi(v) > 0$ for all nonzero $v \in \C$, which is essential for the normalization step.
\end{remark}

%==============================================================================
\section{Proof Strategy: Explicit Construction}
%==============================================================================

The Adviser subagent evaluated four potential approaches:

\begin{enumerate}
    \item \textbf{Direct generalization} from quantum GJWH (viability: 0.5)
    \item \textbf{Convex analysis/separation} approach (viability: 0.7)
    \item \textbf{Duality via $\tmax$ definition} (viability: 0.8)
    \item \textbf{Explicit construction} (viability: 0.85) --- \textbf{Selected}
\end{enumerate}

The explicit construction strategy proceeds as follows:

\begin{enumerate}
    \item Fix a reference measurement setting $x_0$
    \item Normalize the assemblage elements: $\hat{v}_a := v_{a|x_0}/\phi(v_{a|x_0})$
    \item Construct the shared state: $w := \sum_a \phi(v_{a|x_0}) \cdot (\hat{v}_a \otimes \hat{v}_a)$
    \item Define measurements via stochastic rescaling
    \item Verify the representation formula
\end{enumerate}

%==============================================================================
\section{Proof Skeleton (Depth 1 Nodes)}
%==============================================================================

\subsection{Step 1: Normalization (\texttt{:1-setup1})}

\begin{statusbox}[Status: PROPOSED (Challenged)]{yellow}
\textbf{Claim}: Fix $x_0$ and define $\hat{v}_a := \frac{v_{a|x_0}}{\phi(v_{a|x_0})}$ for each $a$ with $\phi(v_{a|x_0}) > 0$, yielding $\hat{v}_a \in \C$ with $\phi(\hat{v}_a) = 1$.

\textbf{Dependencies}: A2, A3, A4

\textbf{Justification}: Definition expansion
\end{statusbox}

\begin{challengebox}
\textbf{Verifier Challenge}: The step silently excludes outcomes where $v_{a|x_0} = 0$ (hence $\phi(v_{a|x_0}) = 0$). The theorem conclusion requires constructing $f_{a|x}$ for ALL outcomes $a$, not just those with nonzero assemblage elements. This is a scope violation.
\end{challengebox}

\textbf{Mathematical Content}: For $v_{a|x_0} \in \C \setminus \{0\}$ and $\phi \in \interior(\Cstar)$, we have $\phi(v_{a|x_0}) > 0$ by the characterization of interior dual functionals. The normalization $\hat{v}_a = v_{a|x_0}/\phi(v_{a|x_0})$ satisfies:
\begin{itemize}
    \item $\hat{v}_a \in \C$ (cones are closed under positive scaling)
    \item $\phi(\hat{v}_a) = \phi(v_{a|x_0})/\phi(v_{a|x_0}) = 1$
\end{itemize}

%------------------------------------------------------------------------------
\subsection{Step 2: Construct $w$ (\texttt{:1-constw})}

\begin{statusbox}[Status: PROPOSED (Challenged)]{yellow}
\textbf{Claim}: Define $w := \sum_a \phi(v_{a|x_0}) \cdot (\hat{v}_a \otimes \hat{v}_a) \in V \otimes V$.

\textbf{Dependencies}: Step 1, A4

\textbf{Justification}: Definition expansion
\end{statusbox}

\begin{challengebox}
\textbf{Verifier Challenge}: The sum index set is ambiguous. For outcomes $a$ where $v_{a|x_0} = 0$, the term $\hat{v}_a$ is undefined. The expression ``$0 \cdot (\text{undefined} \otimes \text{undefined})$'' is not mathematically well-formed.
\end{challengebox}

\textbf{Mathematical Content}: Assuming we restrict to outcomes with $\phi(v_{a|x_0}) > 0$:
\[
w = \sum_{a: \phi(v_{a|x_0}) > 0} \phi(v_{a|x_0}) \cdot (\hat{v}_a \otimes \hat{v}_a)
\]
This is a finite sum (assuming finite outcome set) of rank-1 tensor products weighted by positive scalars.

%------------------------------------------------------------------------------
\subsection{Step 3: Verify $w \in \C \tmax \C$ (\texttt{:1-verw01})}

\begin{statusbox}[Status: PROPOSED (Minor Challenge)]{green}
\textbf{Claim}: $w \in \C \tmax \C$, i.e., for all $f, g \in \Cstar$:
\[
(f \otimes g)(w) = \sum_a \phi(v_{a|x_0}) \cdot f(\hat{v}_a) \cdot g(\hat{v}_a) \geq 0
\]

\textbf{Dependencies}: Step 2, A2, A3

\textbf{Justification}: Algebraic rewrite
\end{statusbox}

\begin{challengebox}
\textbf{Verifier Challenge (Minor)}: The justification \texttt{:algebraic-rewrite} is incorrect. This step invokes the \emph{definition} of the maximal tensor product and properties of the dual cone, not merely algebraic manipulation.
\end{challengebox}

\textbf{Mathematical Content}: This step is \textbf{mathematically correct}. The computation:
\begin{align}
(f \otimes g)(w) &= (f \otimes g)\left(\sum_a \phi(v_{a|x_0}) \cdot (\hat{v}_a \otimes \hat{v}_a)\right) \\
&= \sum_a \phi(v_{a|x_0}) \cdot (f \otimes g)(\hat{v}_a \otimes \hat{v}_a) \\
&= \sum_a \phi(v_{a|x_0}) \cdot f(\hat{v}_a) \cdot g(\hat{v}_a)
\end{align}
Each term is non-negative because:
\begin{itemize}
    \item $\phi(v_{a|x_0}) > 0$ (by construction)
    \item $f(\hat{v}_a) \geq 0$ since $\hat{v}_a \in \C$ and $f \in \Cstar$
    \item $g(\hat{v}_a) \geq 0$ by the same reasoning
\end{itemize}
A sum of non-negative terms is non-negative.

%------------------------------------------------------------------------------
\subsection{Step 4: Construct Measurements (\texttt{:1-constf})}

\begin{rejectedbox}
\textbf{Claim}: For each $(a, x)$, define $f_{a|x} \in V^*$ via the stochastic representation: if $\phi(v_{a|x_0}) > 0$, set
\[
f_{a|x} := \frac{\phi(v_{a|x})}{\phi(v_{a|x_0})} \cdot \phi_{\hat{v}_a}
\]
where $\phi_{\hat{v}_a}(\cdot)$ is a functional ``extending evaluation at $\hat{v}_a$''.

\textbf{Dependencies}: Step 1, A3, A4

\textbf{Justification}: Definition expansion

\textbf{Status}: \textcolor{red}{\textbf{REJECTED}}
\end{rejectedbox}

\begin{challengebox}
\textbf{Verifier Challenge (Critical)}: The notation $\phi_{\hat{v}_a}$ is \textbf{mathematically undefined}. The phrase ``functional extending evaluation at $\hat{v}_a$'' is incoherent:
\begin{itemize}
    \item $\hat{v}_a \in V$ is a vector, not a functional
    \item ``Evaluation at a point'' makes sense for functionals on $V^*$, not on $V$
    \item To interpret $\phi_{\hat{v}_a} \in V^*$, one would need either:
    \begin{enumerate}
        \item An inner product on $V$ (not assumed)
        \item A canonical isomorphism $V \to V^*$ (requires additional structure)
        \item An explicit construction from the given data
    \end{enumerate}
\end{itemize}
Additionally, the case $\phi(v_{a|x_0}) = 0$ is not handled.
\end{challengebox}

%------------------------------------------------------------------------------
\subsection{Step 5: Verify Measurement Properties (\texttt{:1-verfpr})}

\begin{rejectedbox}
\textbf{Claim}: $f_{a|x} \in \Cstar$ for all $a, x$, and $\sum_a f_{a|x} = \phi$ for each $x$.

\textbf{Dependencies}: Step 4, A3, A4

\textbf{Justification}: Algebraic rewrite

\textbf{Status}: \textcolor{red}{\textbf{REJECTED}}
\end{rejectedbox}

\begin{challengebox}
\textbf{Verifier Challenge (Critical)}: This step inherits the undefined notation from Step 4. Moreover, the summation identity $\sum_a f_{a|x} = \phi$ requires:
\[
\sum_a \frac{\phi(v_{a|x})}{\phi(v_{a|x_0})} \cdot \phi_{\hat{v}_a} = \phi
\]
For this to hold for \emph{all} measurement settings $x$, the $\phi_{\hat{v}_a}$ must satisfy very specific constraints. Since the coefficients vary with $x$ but $\phi_{\hat{v}_a}$ do not, this forces $\phi_{\hat{v}_a} = \phi(v_{a|x_0}) \cdot \phi$ (scalar multiples of $\phi$). This is an \textbf{unstated requirement} not established by any dependency.
\end{challengebox}

%------------------------------------------------------------------------------
\subsection{Step 6: Verify Representation (\texttt{:1-reprep})}

\begin{rejectedbox}
\textbf{Claim}: For all $a, x$: $(f_{a|x} \otimes \id)(w) = v_{a|x}$.

\textbf{Dependencies}: Steps 2, 4, 3, A4

\textbf{Justification}: Algebraic rewrite

\textbf{Status}: \textcolor{red}{\textbf{REJECTED --- FUNDAMENTAL ERROR}}
\end{rejectedbox}

\begin{challengebox}
\textbf{Verifier Challenge (FATAL)}: This step contains a \textbf{fundamental structural error}.

Assuming the most charitable interpretation where $\phi_{\hat{v}_a}(\hat{v}_b) = \delta_{ab}$ (Kronecker delta), we compute:
\begin{align}
(f_{a|x} \otimes \id)(w) &= \sum_b \phi(v_{b|x_0}) \cdot f_{a|x}(\hat{v}_b) \cdot \hat{v}_b \\
&= \sum_b \phi(v_{b|x_0}) \cdot \frac{\phi(v_{a|x})}{\phi(v_{a|x_0})} \cdot \delta_{ab} \cdot \hat{v}_b \\
&= \phi(v_{a|x_0}) \cdot \frac{\phi(v_{a|x})}{\phi(v_{a|x_0})} \cdot \hat{v}_a \\
&= \phi(v_{a|x}) \cdot \hat{v}_a \\
&= \phi(v_{a|x}) \cdot \frac{v_{a|x_0}}{\phi(v_{a|x_0})} \\
&= \frac{\phi(v_{a|x})}{\phi(v_{a|x_0})} \cdot v_{a|x_0}
\end{align}

\textbf{This equals $v_{a|x}$ if and only if}:
\[
v_{a|x} = \frac{\phi(v_{a|x})}{\phi(v_{a|x_0})} \cdot v_{a|x_0}
\]
That is, $v_{a|x}$ must be a \textbf{scalar multiple} of $v_{a|x_0}$ for all $x$.

\textbf{This is NOT true for general steering assemblages!} The whole point of steering is that different measurement settings $x$ can produce assemblage elements pointing in genuinely different directions within the cone.
\end{challengebox}

%------------------------------------------------------------------------------
\subsection{QED Step (\texttt{:1-qed001})}

\begin{statusbox}[Status: TAINTED]{red}
\textbf{Claim}: There exist measurements $f_{a|x} \in \Cstar$ with $\sum_a f_{a|x} = \phi$ and $w \in \C \tmax \C$ such that $v_{a|x} = (f_{a|x} \otimes \id)(w)$ for all $a, x$.

\textbf{Dependencies}: Steps 3, 5, 6

\textbf{Status}: TAINTED (depends on rejected steps)
\end{statusbox}

%==============================================================================
\section{Root Cause Analysis}
%==============================================================================

\subsection{The Fundamental Obstruction}

The proof strategy fails because of a fundamental geometric obstruction:

\begin{enumerate}
    \item The construction $w = \sum_a \phi(v_{a|x_0}) \cdot (\hat{v}_a \otimes \hat{v}_a)$ encodes \textbf{only the directional information from setting $x_0$}.

    \item The measurements $f_{a|x}$ can only provide \textbf{scalar factors} $\phi(v_{a|x})/\phi(v_{a|x_0})$ that rescale the output.

    \item The map $f \mapsto (f \otimes \id)(w)$ is \textbf{linear} in $f$, so different measurements can only move along ``lines'' in the output space.

    \item But genuine steering assemblages have $v_{a|x}$ pointing in \textbf{genuinely different directions} for different $x$---not just scalar multiples of a fixed direction.
\end{enumerate}

\subsection{Why Quantum GJWH Works}

In the quantum case (PSD cone), the proof works because:

\begin{enumerate}
    \item A \textbf{purification} $|\Psi\rangle$ of the reduced state $\rho_B$ exists.

    \item By the Schr\"odinger--HJW theorem, \emph{any} ensemble decomposition of $\rho_B$ can be realized by some measurement on the purifying system.

    \item The key property: the \textbf{same} purification works for all measurement settings simultaneously.

    \item This relies on the special structure of PSD cones: self-duality, homogeneity, and the rich ``entanglement'' structure of the tensor product.
\end{enumerate}

\subsection{What Would Be Needed}

For the theorem to hold for general proper cones, one would need:

\begin{enumerate}
    \item A ``purification'' or ``extension'' concept for proper cones
    \item A proof that $\C \tmax \C$ contains enough ``entangled'' elements
    \item A non-constructive existence argument (separation theorems)
\end{enumerate}

The explicit construction approach \textbf{fundamentally cannot work} because it builds $w$ from only one measurement setting's data.

%==============================================================================
\section{Verification Summary}
%==============================================================================

\begin{center}
\begin{tabular}{llll}
\toprule
\textbf{Node ID} & \textbf{Description} & \textbf{Status} & \textbf{Issue Severity} \\
\midrule
\texttt{:0-a00001} & $V$ finite-dimensional & Verified & --- \\
\texttt{:0-a00002} & $\C$ proper cone & Verified & --- \\
\texttt{:0-a00003} & $\phi \in \interior(\Cstar)$ & Verified & --- \\
\texttt{:0-a00004} & Steering assemblage & Verified & --- \\
\texttt{:1-setup1} & Normalization & Proposed & Moderate \\
\texttt{:1-constw} & Construct $w$ & Proposed & Moderate \\
\texttt{:1-verw01} & Verify $w \in \C \tmax \C$ & Proposed & Minor \\
\texttt{:1-constf} & Construct measurements & \textcolor{red}{Rejected} & Critical \\
\texttt{:1-verfpr} & Verify measurement properties & \textcolor{red}{Rejected} & Critical \\
\texttt{:1-reprep} & Verify representation & \textcolor{red}{Rejected} & \textbf{Fatal} \\
\texttt{:1-qed001} & QED & Tainted & --- \\
\bottomrule
\end{tabular}
\end{center}

%==============================================================================
\section{Conclusion}
%==============================================================================

The explicit construction proof strategy for \cref{thm:main} fails due to a fundamental structural obstruction. The construction of $w$ from a single measurement setting's data cannot encode the full ``rotational'' structure of a general steering assemblage.

This finding suggests that \cref{thm:main} is likely \textbf{false} for general proper cones, though a rigorous counterexample is needed to definitively settle the question. The theorem \emph{is} true for the special case of PSD cones (quantum mechanics) due to the rich entanglement structure of matrix tensor products.

\vspace{1em}
\begin{center}
\textit{Detection beats sycophancy. Finding an error is a success, not a failure.}\\
--- Alethfeld Proof Orchestrator v5.1
\end{center}

\end{document}
